\chapter{How to add formatted code}

\section{Basic Example}
This is some random MATLAB code. You can reference it just like "Listing~\ref{lst:matlabCode}" 
or "Listing~\ref{lst:inlineExample}".


\subsection*{Example 1: Loading from External File}
The following code is loaded from an external MATLAB file, showing only lines 7-20:

\lstinputlisting[caption={aMatlabFile.m}, 
                 label={lst:matlabCode}, 
                 linerange={7-20}, 
                 firstnumber=7]{src/code/aMatlabFile.m}

\subsection*{Example 2: Inline Code Block}
You can also write code directly in your document:
\begin{lstlisting}[
    caption = {Simple MATLAB function example},
    label   = {lst:inlineExample},
    language = MATLAB
]
function y = myFunction(x)
    % MYFUNCTION Computes square of input
    % Input:  x - numeric value
    % Output: y - x squared
    
    y = x^2;
    
    if y > 100
        disp('Large value detected!');
    end
end
\end{lstlisting}

\section{Creating Custom Formats}
            
You can create your own custom formats by defining new styles. Comprehensive documentation and examples for the listings package can be found at: \url{https://ctan.org/pkg/listings}.
